\documentclass{article}

% NeurIPS 2026 Education Track. The "education" option is non-anonymous,
% which matches the track's single-blind review with author names included.
\usepackage[education, final]{neurips_2026}

\usepackage[utf8]{inputenc}
\usepackage[T1]{fontenc}
\usepackage{hyperref}
\usepackage{url}
\usepackage{amsfonts}

\title{Right Answer, Wrong Concept:\\Teaching Reasoning Shortcuts in Neurosymbolic Learning}

\author{%
  Victoria McCray\\
  Northeastern University\\
  \texttt{mccray.v@northeastern.edu}
}

\begin{document}

\maketitle

\section{The concept}

Reasoning shortcuts are high-performing solutions in neurosymbolic learning in
which learned concepts acquire unintended semantics while remaining compatible
with task supervision and symbolic reasoning. Following Marconato et al.
(2023), a reasoning shortcut is a learned concept distribution satisfying two
conditions jointly: it attains maximal or near maximal likelihood on the
training objective, and it diverges from ground truth concept semantics. Both
conditions matter, because a low concept score alone is equally consistent with
an undertrained or poorly optimised model, an ordinary engineering failure with
an ordinary fix. Separating those two cases is the central skill taught here.

The resource uses progressively richer interactive examples, from exhaustive
Boolean reasoning to MNIST and simulated driving, to teach how such shortcuts
arise, how to distinguish them from optimisation failure, and how to evaluate
and mitigate them. The resource is organised around four factors identified by
Marconato et al.: prior knowledge, architecture, learning objective, and data
support, each illustrated through a worked example. Throughout, task accuracy
stays near its ceiling while concept quality ranges from 0.998 to 0.000, so
learners see one metric hide four mechanisms needing four different diagnostics.

\section{Leveling and prerequisite knowledge}

The materials assume basic familiarity with supervised machine learning, neural
networks, classification, train and test evaluation, and metrics such as
accuracy and F1 score, together with elementary probability and Boolean logic.
Prior knowledge of neurosymbolic learning, logic programming, or reasoning
shortcuts is not required, and each term is defined at the point of use. The
setup is introduced progressively through visual explanations and executable
examples before moving to diagnosis and mitigation. The conceptual portions are
accessible to a broad technical audience, while the notebook exercises and
architecture comparisons are best suited to advanced undergraduate and early
graduate instruction, as well as early-career scientists.

\section{Learning objectives and outcomes}

By the end of the material a learner will be able to:

\begin{itemize}
  \item \textbf{Remember.} State the two-condition definition of a reasoning
    shortcut and name the four factors governing whether one exists.
  \item \textbf{Understand.} Explain reasoning shortcuts, distinguish task
    performance from concept quality, and connect both to construct validity as
    understood in measurement and in the fairness literature.
  \item \textbf{Apply.} Run the diagnostics on a trained model, read a concept
    confusion matrix, and experiment with mitigations such as concept
    supervision and multi-task learning.
  \item \textbf{Analyse.} Investigate why additional data may not eliminate a
    shortcut when concept support is already exhaustive, and examine the
    implications for out-of-distribution generalisation.
  \item \textbf{Evaluate.} Judge whether a result establishes concept grounding:
    ruling out optimisation failure with a ceiling check, recognising that
    constraint satisfaction can hold by construction, and rejecting averaged
    scores that hide a concept at zero.
  \item \textbf{Create.} Design a task and rule inducing a shortcut by
    construction, and specify a diagnostic protocol naming the factor at issue
    and the measurement exposing it.
\end{itemize}

\section{Linked papers}

\begin{itemize}
  \item E. Marconato, S. Teso, A. Vergari, A. Passerini. \emph{Not All
    Neuro-Symbolic Concepts Are Created Equal.} NeurIPS 2023, arXiv:2305.19951.
    Formalises reasoning shortcuts and identifies the four governing factors.
  \item S. Bortolotti, E. Marconato, T. Carraro, P. Morettin, et al. \emph{A
    Neuro-Symbolic Benchmark Suite for Concept Quality and Reasoning Shortcuts.}
    NeurIPS Datasets and Benchmarks 2024, arXiv:2406.10368. Contributes
    \texttt{rsbench}, which counts the shortcuts a task admits.
  \item E. Marconato, S. Bortolotti, E. van Krieken, A. Vergari, A. Passerini,
    S. Teso. \emph{BEARS Make Neuro-Symbolic Models Aware of their Reasoning
    Shortcuts.} UAI 2024, arXiv:2402.12240. Calibrates rather than eliminates.
  \item X. Yang et al. \emph{Analysis for Abductive Learning and Neural-Symbolic
    Reasoning Shortcuts.} ICML 2024. Bounds shortcut risk by knowledge base
    complexity, sample size and hypothesis space.
  \item E. van Krieken, P. Minervini, E. Ponti, A. Vergari. \emph{Neurosymbolic
    Reasoning Shortcuts under the Independence Assumption.} NeSy 2025,
    arXiv:2507.11357. Proves independence blocks uncertainty over combinations.
  \item A. Takemura, K. Inoue, M. Nishino. \emph{Constraint-Based Analysis of
    Reasoning Shortcuts in Neurosymbolic Learning.} 2026, arXiv:2604.23377.
    Deciding shortcut freeness is coNP-complete.
\end{itemize}

\section{Teaching materials summary}

All materials are at
\url{https://github.com/victoriamccray/neurosymbolic-reasoning-shortcuts}
under the MIT licence.

\textbf{Interactive lesson.} A 31 slide reveal.js deck in eight sections,
opening on Clever Hans as a cross-domain hook and closing on a learner
challenge. It runs offline with no build step, carries alternative text on every
figure, respects \texttt{prefers-reduced-motion}, and measures each slide so
none can overflow when projected.

\textbf{Notebook.} A 26 cell Jupyter notebook making every diagnostic runnable
end to end: the XOR enumeration, the entangled against disentangled encoder
comparison with per concept F1 and confusion matrices, and the confounded
driving scenario with a shift test and an interactive supervision sweep. Each
cell is paired with a plain language explanation of why it is there.

\textbf{Evaluation guidance.} Mitigation is a measured trade-off, not a monotone
improvement: concept supervision raises out-of-distribution accuracy from 0.535
to 0.855 while in-distribution accuracy falls from 0.988 to 0.920. A closing
checklist separates evidence supporting a semantic claim from evidence leaving
it open, and one result is deliberately marked as cited rather than reproduced.

\end{document}
